\subsection{Restricciones}

\subsubsection{Regla: prohibido capturar o lanzar \texttt{Exception} genérica}

\textbf{Regla:} Prohibido capturar o lanzar \texttt{System.Exception} u otro tipo genérico de excepción base, en firma o en \texttt{catch}; siempre el tipo más específico.

\textbf{Descripción:} Regla de análisis \texttt{CA1031} de .NET: las excepciones genéricas ocultan problemas en tiempo de ejecución y dificultan la depuración (Microsoft, 2016).

\textbf{Código correcto:}
\begin{lstlisting}[style=csharpstyle]
public void EstablishConnection()
{
    try
    {
        _database.Connect();
    }
    catch (TimeoutException ex)
    {
        throw new GameStateException("Connection timed out.", ex);
    }
}
\end{lstlisting}

\textbf{Código incorrecto:}
\begin{lstlisting}[style=csharpstyle]
public void EstablishConnection()
{
    try
    {
        _database.Connect();
    }
    catch (Exception ex)
    {
        throw new GameStateException("Connection failed.", ex);
    }
}
\end{lstlisting}

\subsubsection{Regla: sufijo \texttt{Exception} y tres constructores obligatorios}

\textbf{Regla:} Toda excepción personalizada termina en \texttt{Exception} e incluye tres constructores: sin parámetros, con mensaje, y con mensaje más excepción interna (Microsoft, 2025a).

\textbf{Código correcto:}
\begin{lstlisting}[style=csharpstyle]
public sealed class GameStateException : Exception
{
    public GameStateException()
    {
    }

    public GameStateException(string message) : base(message)
    {
    }

    public GameStateException(string message, Exception innerException) : base(message, innerException)
    {
    }
}
\end{lstlisting}

\textbf{Código incorrecto:}
\begin{lstlisting}[style=csharpstyle]
public sealed class GameStateError : Exception
{
    public GameStateError(string message) : base(message)
    {
    }
}
\end{lstlisting}
